\documentclass[mode=notes, palette=ocean, lang=german, fontsize=11pt]{modernclassnotes}

\course{LaTeX-Paketdokumentation}
\coursecode{MCN v1.2.0}
\professor{Modern Class Notes Entwickler}
\institution{Open Source TeX Initiative}
\term{Version 2026}
\docsubtitle{Benutzerhandbuch \& Architektur-Leitfaden für das Class-Notes-Paket}
\title{Das \texttt{modernclassnotes} Paket \& Dokumentenklasse}

\begin{document}

\maketitle

\tableofcontents
\newpage

\section{Einführung}

Das \texttt{modernclassnotes}-Paket und die zugehörige Dokumentenklasse bieten ein modernes, hochästhetisches und quelloffenes LaTeX-Framework, das speziell für Professoren, Dozenten und Lehrkräfte entwickelt wurde. Es ermöglicht die mühelose Erstellung und Verteilung von Vorlesungsskripten, Handouts und Übungsblättern.

\begin{takeaways}
Wichtigste Merkmale von \texttt{modernclassnotes}:
\begin{itemize}
    \item \textbf{Vier spezialisierte Modi}: \texttt{notes} (vollständige Vorlesungsskripte), \texttt{handout} (1--2-seitige Zusammenfassungen), \texttt{worksheet} (Übungsblätter mit Lösungsschlüsseln) und \texttt{chapternotes} (kapitelbasierte Zusammenfassungen).
    \item \textbf{Kuratierte Farbpaletten}: Wechseln Sie zwischen 8 Farbthemen (\texttt{ocean}, \texttt{nord}, \texttt{midnight}, \texttt{emerald}, \texttt{burgundy}, \texttt{amethyst}, \texttt{purple}, \texttt{mono}) mit einer einzigen Option.
    \item \textbf{Mehrsprachige Lokalisierung}: Nativ unterstützte Sprachen mit \texttt{lang=english}, \texttt{lang=german}, \texttt{lang=french} und \texttt{lang=chinese}.
    \item \textbf{Keine komplexen Abhängigkeiten}: Native TeX/LaTeX-Box-Engine für maximale Portabilität unter TeX Live, MiKTeX und Overleaf.
    \item \textbf{Vorlesungs- \& Kapitel-Verfolgung}: Automatische Vorlesungs- und Kapitelbanner mit optionalen Lernzielen via \texttt{\textbackslash lecture}, \texttt{\textbackslash chapter} und \texttt{\textbackslash subchapter}.
\end{itemize}
\end{takeaways}

\section{Klassenoptionen \& Konfiguration}

Laden Sie die Klasse im Dokumentenkopf:
\begin{codebox}[Einbinden der Klasse]
\begin{lstlisting}[language=TeX]
\documentclass[mode=notes, palette=ocean, lang=german, fontsize=11pt]{modernclassnotes}
\end{lstlisting}
\end{codebox}

\subsection{Verfügbare Optionen}
\begin{table}[h]
\centering
\begin{tabular}{l|p{9.5cm}|l}
\toprule
\textbf{Schlüssel} & \textbf{Werte} & \textbf{Standard} \\
\midrule
\texttt{mode} & \texttt{notes}, \texttt{handout}, \texttt{worksheet}, \texttt{chapternotes} & \texttt{notes} \\ \hline
\texttt{palette} & \texttt{ocean}, \texttt{midnight}, \texttt{nord}, \texttt{emerald}, \texttt{burgundy}, \texttt{amethyst}, \texttt{purple}, \texttt{mono} & \texttt{ocean} \\ \hline
\texttt{fontsize} & \texttt{10pt}, \texttt{11pt}, \texttt{12pt} & \texttt{11pt} \\ \hline
\texttt{font} & \texttt{default}, \texttt{palatino}, \texttt{charter}, \texttt{sans}/\texttt{helvet}, \texttt{kpfonts}, \texttt{pagella}, \texttt{sourceserif}, \texttt{fira}, \texttt{libertine}, \texttt{utopia} & \texttt{default} \\ \hline
\texttt{lang} & \texttt{english} (\texttt{en}), \texttt{german} (\texttt{de}), \texttt{french} (\texttt{fr}), \texttt{chinese} (\texttt{zh}) & \texttt{english} \\ \hline
\texttt{official} & \texttt{true}, \texttt{false} & \texttt{false} \\ \hline
\texttt{showsolutions} & \texttt{true}, \texttt{false} & \texttt{true} \\ \hline
\texttt{parskip} & \texttt{true}, \texttt{false} & \texttt{true} \\
\bottomrule
\end{tabular}
\caption{Klassenoptionen und Parameterwerte}
\end{table}

\section{Kurs-Metadaten Befehle}

Geben Sie Ihre Kursdaten in der Präambel an:
\begin{codebox}[Metadaten-Einrichtung]
\begin{lstlisting}[language=TeX]
\course{Theoretische Informatik}
\coursecode{CS 102}
\professor{Prof. Dr. Markus Weber}
\institution{Technische Universität München}
\term{Wintersemester 2026}
\docsubtitle{Vorlesungsskript \& Übungsaufgaben}
\title{Einführung in die Algorithmik}
\end{lstlisting}
\end{codebox}

\section{Vorlesungs- \& Kapitel-Systeme}

Verwenden Sie den Befehl \texttt{\textbackslash lecture}, um das Skript in Vorlesungseinheiten zu gliedern:
\begin{codebox}[Vorlesungs-Makro]
\begin{lstlisting}[language=TeX]
\lecture[15. Oktober 2026]{1}{Komplexitätstheorie \& O-Notation}
\end{lstlisting}
\end{codebox}

Für kapitelbasierte Skripte (\texttt{mode=chapternotes}) nutzen Sie \texttt{\textbackslash chapter} und \texttt{\textbackslash subchapter}:
\begin{codebox}[Kapitel-Makro mit Lernzielen]
\begin{lstlisting}[language=TeX]
\chapter{1}{Digitale Systeme und Binärzahlen}{%
  Binärzahlen verstehen.,%
  Umwandlung zwischen Zahlenbasis-Systemen.,%
  Bildung von Komplementen.%
}

\subchapter{2}{Binärzahlen}
\end{lstlisting}
\end{codebox}

\section{Pädagogische Umgebungen}

\begin{definition}[Mathematische \& Didaktische Umgebungen]
\texttt{modernclassnotes} bietet strukturierte Umgebungen für lehrreiche Inhalte:
\begin{itemize}
    \item \texttt{theorem} (Satz), \texttt{lemma}, \texttt{proposition}, \texttt{corollary} (Korollar), \texttt{proof} (Beweis)
    \item \texttt{definition}, \texttt{example} (Beispiel), \texttt{remark} (Bemerkung), \texttt{warning} (Warnung), \texttt{tip} (Tipp), \texttt{takeaways}
    \item \texttt{exercise} (Übung), \texttt{solution} (Lösung), \texttt{codebox} (Quellcode)
\end{itemize}
\end{definition}

\begin{example}[Beispiel für Satz-Umgebung]
\begin{lstlisting}[language=TeX]
\begin{theorem}[Eulersche Identität]
Für jede reelle Zahl $\theta$ gilt $e^{i\theta} = \cos\theta + i\sin\theta$.
\end{theorem}
\end{lstlisting}
\end{example}

\section{Open-Source Richtlinien}

Zur Weitergabe und Bereitstellung dieses Pakets:
\begin{enumerate}
    \item Erhalten Sie die LPPL 1.3c / MIT Lizenzbedingungen in \texttt{LICENSE}.
    \item Stellen Sie kompilierte Beispiel-PDFs in \texttt{examples/} für Vorschauen bereit.
\end{enumerate}

\end{document}
